La \texttt{ExamInstance} atraviesa una secuencia de estados que refleja
su progreso desde la entrega física hasta el archivo definitivo del
examen. La \cref{fig:scdee-exam-state} resume la máquina de estados
completa, incluidas las transiciones de retroceso que el sistema permite
(recorrección de un examen ya calificado, reapertura de la revisión sobre
un examen finalizado).

\begin{figure}[Máquina de estados de la \texttt{ExamInstance}]%
              {fig:scdee-exam-state}%
              {Máquina de estados de la \texttt{ExamInstance}.}
  \image{}{0.85\textheight}{mermaid/exam_state}
\end{figure}

Las transiciones se validan contra una tabla central de transiciones
legales antes de aplicarse, lo que impide saltos arbitrarios entre
estados y centraliza la lógica del ciclo en un único punto.

A nivel de página, cada \texttt{ExamPage} mantiene su propio estado más
sencillo (\texttt{PENDING\_RECOGNITION}, \texttt{RECOGNIZED},
\texttt{ORPHAN}, \texttt{DISCARDED}), independiente del estado de la
instancia a la que pertenece.

Las instancias y las páginas mantienen, además del estado, un conjunto
de incidencias declarables (\texttt{MISSING\_PAGE},
\texttt{STUDENT\_NOT\_IDENTIFIED}, \texttt{DUPLICATE\_STUDENT} a nivel
de instancia; \texttt{QR\_READ\_ERROR}, \texttt{QR\_MISMATCH},
\texttt{ORPHAN\_PAGE}, \texttt{EXTRA\_PAGE} a nivel de página). Estas
incidencias son ortogonales al estado, pero interactúan con el ciclo:
una instancia con incidencias activas no puede transitar a
\texttt{PUBLISHED}. La resolución manual de cada incidencia por parte
del coordinador (asignación del alumno, descarte de páginas extra,
edición del modelo asignado) limpia el correspondiente \textit{flag}
y desbloquea la transición.
